% Mapeos de caracteres unicode que las fuentes TeX Gyre no cubren
% (se usa junto con tema_bugambilia.tex en los decks de este curso).
\usepackage{newunicodechar}
\newunicodechar{→}{\ensuremath{\rightarrow}}
\newunicodechar{←}{\ensuremath{\leftarrow}}
\newunicodechar{↔}{\ensuremath{\leftrightarrow}}
\newunicodechar{⇒}{\ensuremath{\Rightarrow}}
\newunicodechar{ℓ}{\ensuremath{\ell}}
\newunicodechar{ϕ}{\ensuremath{\phi}}
\newunicodechar{φ}{\ensuremath{\varphi}}
\newunicodechar{θ}{\ensuremath{\theta}}
\newunicodechar{α}{\ensuremath{\alpha}}
\newunicodechar{β}{\ensuremath{\beta}}
\newunicodechar{γ}{\ensuremath{\gamma}}
\newunicodechar{δ}{\ensuremath{\delta}}
\newunicodechar{Δ}{\ensuremath{\Delta}}
\newunicodechar{λ}{\ensuremath{\lambda}}
\newunicodechar{μ}{\ensuremath{\mu}}
\newunicodechar{π}{\ensuremath{\pi}}
\newunicodechar{σ}{\ensuremath{\sigma}}
\newunicodechar{τ}{\ensuremath{\tau}}
\newunicodechar{×}{\ensuremath{\times}}
\newunicodechar{·}{\ensuremath{\cdot}}
\newunicodechar{∈}{\ensuremath{\in}}
\newunicodechar{∂}{\ensuremath{\partial}}
\newunicodechar{−}{\ensuremath{-}}
\newunicodechar{±}{\ensuremath{\pm}}
\newunicodechar{≈}{\ensuremath{\approx}}
\newunicodechar{≠}{\ensuremath{\neq}}
\newunicodechar{≤}{\ensuremath{\leq}}
\newunicodechar{≥}{\ensuremath{\geq}}
\newunicodechar{∞}{\ensuremath{\infty}}
\newunicodechar{‖}{\ensuremath{\|}}
\newunicodechar{‑}{\mbox{-}}
\newunicodechar{✓}{\textcolor{green!55!black}{\ensuremath{\checkmark}}}
\newunicodechar{✅}{\textcolor{green!55!black}{\ensuremath{\checkmark}}}
\newunicodechar{❌}{\textcolor{red!75!black}{\ensuremath{\times}}}

% Salvaguarda global: ninguna imagen crece más allá del alto del frame
\usepackage{graphicx}
\setkeys{Gin}{keepaspectratio, totalheight=0.72\textheight}

% Cap duro: toda imagen (markdown o knitr) se ajusta también al alto del frame.
% Se re-aplica keepaspectratio al final para que width explícito + cap convivan.
\makeatletter
\let\ucompat@incgfx\includegraphics
\renewcommand{\includegraphics}[2][]{%
  {\centering\ucompat@incgfx[#1,keepaspectratio,totalheight=0.62\textheight]{#2}\par}}
\makeatother
